\section{The Playbook: Levers by Position and Setup}\label{sec:playbook}
\IfFileExists{part2/tables/playbook_macros.tex}{% Generated by code/playbook/run_archetypes.py. Do not edit by hand.
\newcommand{\SpendRatioPB}{0.99}
\newcommand{\SpendRatioPBAll}{3.3}
\newcommand{\pbAFiveGap}{research and development}
\newcommand{\pbAFiveGi}{\ensuremath{-}4.2}
\newcommand{\pbAFivePubCaph}{0.27}
\newcommand{\pbAFourGA}{19}
\newcommand{\pbAFourSigma}{18}
\newcommand{\pbAOneDur}{0.62}
\newcommand{\pbAOneGA}{24}
\newcommand{\pbAOneSM}{50}
\newcommand{\pbAOneSigma}{25}
\newcommand{\pbASevenAOL}{\ensuremath{-}1.53}
\newcommand{\pbASevenAOLprice}{2.30}
\newcommand{\pbASevenGi}{\ensuremath{-}5.1}
\newcommand{\pbASevenGorg}{3}
\newcommand{\pbASevenGrep}{126}
\newcommand{\pbASevenReorg}{6.0}
\newcommand{\pbASevenRorg}{37}
\newcommand{\pbASevenRrep}{160}
\newcommand{\pbASevenXpc}{0.77}
\newcommand{\pbASixGi}{36.3}
\newcommand{\pbASixPubCaph}{0.61}
\newcommand{\pbASixXPE}{11.3}
\newcommand{\pbASixXVC}{17.6}
\newcommand{\pbASixXhat}{3.3}
\newcommand{\pbAThreeGap}{research and development}
\newcommand{\pbAugStepsLow}{23}
\newcommand{\pbAutoStepsHigh}{15}
\newcommand{\pbAutoStepsHighShare}{2}
\newcommand{\pbBandwidthUseMax}{1.19}
\newcommand{\pbBandwidthUseMin}{1.01}
\newcommand{\pbCapShareLow}{43}
\newcommand{\pbCovDWA}{207}
\newcommand{\pbCovDWAAll}{2{,}087}
\newcommand{\pbCovSteps}{136}
\newcommand{\pbCovStepsPop}{130}
\newcommand{\pbCovUnits}{662}
\newcommand{\pbCutGmed}{12.7}
\newcommand{\pbCutGseventyfive}{21.7}
\newcommand{\pbCutGtwentyfive}{5.4}
\newcommand{\pbCutMmed}{2.5}
\newcommand{\pbCutMseventyfive}{11.7}
\newcommand{\pbCutMtwentyfive}{\ensuremath{-}10.3}
\newcommand{\pbCutNEtwo}{365}
\newcommand{\pbCutNpassforty}{10}
\newcommand{\pbCutNpop}{435}
\newcommand{\pbCutRninetyfive}{44.6}
\newcommand{\pbCutRseventyfive}{27.7}
\newcommand{\pbDurAll}{yes}
\newcommand{\pbGEAheadNetList}{A1}
\newcommand{\pbGEFirstFamList}{A1}
\newcommand{\pbICFRAugHigh}{14}
\newcommand{\pbICFRAutoHigh}{2}
\newcommand{\pbICFRHouseHigh}{21}
\newcommand{\pbICFROutHigh}{54}
\newcommand{\pbICFRReadShare}{31}
\newcommand{\pbICFRRelShare}{45}
\newcommand{\pbICFRRelShareNow}{0}
\newcommand{\pbInvFiveHighBind}{reliability: detection}
\newcommand{\pbInvFiveHighDnow}{0.65}
\newcommand{\pbInvFiveHighDstar}{0.89}
\newcommand{\pbInvFiveHighFloor}{0.50}
\newcommand{\pbInvFiveHighMode}{O}
\newcommand{\pbInvFiveHighRead}{0.83}
\newcommand{\pbInvFiveHighResid}{17.0}
\newcommand{\pbInvFiveHighTier}{B}
\newcommand{\pbInvFiveHighVol}{35{,}411}
\newcommand{\pbInvFiveLowBind}{readiness}
\newcommand{\pbInvFiveLowDnow}{0.65}
\newcommand{\pbInvFiveLowDstar}{0.92}
\newcommand{\pbInvFiveLowFloor}{0.50}
\newcommand{\pbInvFiveLowMode}{O}
\newcommand{\pbInvFiveLowRead}{0.37}
\newcommand{\pbInvFiveLowResid}{17.0}
\newcommand{\pbInvFiveLowTier}{B}
\newcommand{\pbInvFiveLowVol}{35{,}411}
\newcommand{\pbInvGate}{3}
\newcommand{\pbInvTwoHighBind}{reliability: detection}
\newcommand{\pbInvTwoHighDnow}{0.61}
\newcommand{\pbInvTwoHighDstar}{0.85}
\newcommand{\pbInvTwoHighFloor}{0.50}
\newcommand{\pbInvTwoHighMode}{O}
\newcommand{\pbInvTwoHighRead}{0.83}
\newcommand{\pbInvTwoHighResid}{12.6}
\newcommand{\pbInvTwoHighTier}{B}
\newcommand{\pbInvTwoHighVol}{268{,}428}
\newcommand{\pbInvTwoLowBind}{readiness}
\newcommand{\pbInvTwoLowDnow}{0.65}
\newcommand{\pbInvTwoLowDstar}{0.90}
\newcommand{\pbInvTwoLowFloor}{0.50}
\newcommand{\pbInvTwoLowMode}{O}
\newcommand{\pbInvTwoLowRead}{0.37}
\newcommand{\pbInvTwoLowResid}{12.6}
\newcommand{\pbInvTwoLowTier}{B}
\newcommand{\pbInvTwoLowVol}{268{,}428}
\newcommand{\pbKeepMax}{0.63}
\newcommand{\pbKeepMin}{0.49}
\newcommand{\pbMaintBase}{25.1}
\newcommand{\pbNBindBandwidth}{six}
\newcommand{\pbNCoRTopTwo}{none}
\newcommand{\pbNDefinedFam}{six}
\newcommand{\pbNDefinedHyb}{six}
\newcommand{\pbNDefinedPC}{five}
\newcommand{\pbNDefinedPE}{six}
\newcommand{\pbNDefinedPub}{six}
\newcommand{\pbNDefinedVC}{six}
\newcommand{\pbNEngFirstRevAll}{one}
\newcommand{\pbNEngFirstRevNet}{six}
\newcommand{\pbNEngFirstRevNone}{six}
\newcommand{\pbNGAAheadGENet}{five}
\newcommand{\pbNGAFirstRevAll}{none}
\newcommand{\pbNGAFirstRevNet}{none}
\newcommand{\pbNGAFirstRevNone}{none}
\newcommand{\pbNGEAheadGAAll}{six}
\newcommand{\pbNGEAheadGANet}{one}
\newcommand{\pbNGEAheadGANone}{none}
\newcommand{\pbNGEFirstFam}{one}
\newcommand{\pbNGEFirstHyb}{none}
\newcommand{\pbNGEFirstPC}{none}
\newcommand{\pbNGEFirstPE}{none}
\newcommand{\pbNGEFirstPub}{none}
\newcommand{\pbNGEFirstRevAll}{five}
\newcommand{\pbNGEFirstRevNet}{none}
\newcommand{\pbNGEFirstRevNone}{none}
\newcommand{\pbNGEFirstRevealed}{none}
\newcommand{\pbNGEFirstVC}{none}
\newcommand{\pbNGEThenEng}{none}
\newcommand{\pbNPENonePay}{three}
\newcommand{\pbNSMAutoAheadGA}{six}
\newcommand{\pbNSMMaintAboveGA}{six}
\newcommand{\pbPCRate}{22.3}
\newcommand{\pbPCUndefined}{A6}
\newcommand{\pbPENonePay}{A1, A4, A5}
\newcommand{\pbPERate}{19.9}
\newcommand{\pbPtElasMax}{0.51}
\newcommand{\pbPtElasMin}{0.37}
\newcommand{\pbPtRate}{0.50}
\newcommand{\pbReadShareLow}{31}
\newcommand{\pbRelShareHigh}{42}
\newcommand{\pbSMNetZeroList}{A3, A6}
\newcommand{\pbSpendRatioNet}{0.99}
\newcommand{\pbStepsPhysicalAll}{27}
\newcommand{\pbStepsPhysicalFinance}{10}
\newcommand{\pbStepsPhysicalFinanceList}{P01.4, P01.5, P02.6, P03.3, P04.2, P04.4, P06.3, P19.1, P19.2, P19.3}
\newcommand{\pbStepsPhysicalFinanceSome}{P02.5}
\newcommand{\pbTierStepA}{51}
\newcommand{\pbTierStepB}{24}
\newcommand{\pbTierStepC}{26}
\newcommand{\pbTierUnitA}{55}
\newcommand{\pbTierUnitB}{12}
\newcommand{\pbTierUnitC}{33}
\newcommand{\pbVCRate}{27.0}
\newcommand{\pbXFamMax}{22.4}
\newcommand{\pbXFamMin}{4.3}
\newcommand{\pbXPCMax}{5.4}
\newcommand{\pbXPCMin}{0.1}
\newcommand{\pbXPEMax}{11.3}
\newcommand{\pbXPEMin}{\ensuremath{-}0.1}
\newcommand{\pbXVCMax}{17.6}
\newcommand{\pbXVCMin}{0.9}
\newcommand{\pbXdagMax}{13.1}
\newcommand{\pbXdagMin}{3.3}
\newcommand{\pbXhatMax}{4.5}
\newcommand{\pbXhatMin}{3.3}
}{}% generated by code/playbook/run_archetypes.py

The preceding sections value each lever of an AI transformation on its own. This section asks how a strategy should differ between firms at different points of Part~I's growth--margin plane and with different organizations, and which levers pay, and which do not, where. A firm's state has two parts. Its \emph{position} is its point $(g,m)$ together with the Part~I parameters that rationalize it, and its \emph{setup} is the set of organizational dimensions treated in Sections~\ref{sec:tasks} to \ref{sec:acquirer} (Section~\ref{sec:playbook-state}). The position ranks the levers by marginal value through Part~I's envelope results, and the ranking switches at thresholds in the owner's growth weight, the remaining headroom, the durable share of a saving and the growth rate (Proposition~\ref{prop:playbook-rank}). The setup then removes levers through gates that hold whatever the position (Proposition~\ref{prop:playbook-gates} and Corollary~\ref{cor:playbook-nopay}). Archetypes defined at quantiles of Part~I's SEC panel make the results concrete (Section~\ref{sec:playbook-archetypes}), the engine generates the playbook table (Section~\ref{sec:playbook-table}), and an executable record for every process step and activity carries each recommendation down to the level at which work is assigned (Section~\ref{sec:playbook-steps}). Four conclusions stand out. Whether automating the growth engine or cutting overhead comes first depends on how much selling and marketing buys growth: on Part~I's central basis, which counts as growth spend only selling and marketing above a maintenance level, cutting general and administrative cost comes first at \pbNGAAheadGENet{} of the six panel positions, and automating the growth engine leads only where selling and marketing is well above that level (engineering, valued as if all R\&D were maintenance, ranks above both). Given the basis, the ranking is a statement about the owner's growth weight relative to the weight that rationalizes the firm's own spending, so the same firm can rank its levers differently under a public owner, a sponsor or a permanent acquirer. Savings in the growth engine are transient and savings in overhead are durable, so the ranking also turns as headroom is used up. Cost-of-revenue automation keeps only the part of a saving that customers do not take through prices, so it is among the first two levers at \pbNCoRTopTwo{} of the panel positions. And most step-level decisions are set by setup gates, feasibility, reliability, readiness and volume, rather than by position. Proofs are in Appendix~\ref{app:proofs-playbook}; \texttt{code/verify\_part2\_playbook.py} checks every closed form, threshold and generated number, and \texttt{code/playbook/run\_archetypes.py} generates every number in this section through the engine's public interface.

\subsection{The state of the firm}\label{sec:playbook-state}

\begin{definition}[State]\label{def:playbook-state}
The \emph{position} of a firm is its revenue growth $g$ and free-cash-flow margin $m$ after stock-based compensation, the parameters $\vartheta=(m_0,k,\ln\lambda,G_i)$ of Assumption~\ref{ass:program} with the spend elasticity $\alpha$, its growth spend $\sigma$, its cost lines (general and administrative cost $a_{\rm GA}$, maintenance R\&D $a_{\rm RD}$ and maintenance selling and marketing $a_{\rm SM}$ in $a$, human service cost $c_h$ in cost of revenue), and two weights: the weight $\hat X$ that rationalizes its spending and the investor's screening weight $X^{\rm scr}$ of Lemma~\ref{lem:valuation2}. The \emph{setup} is the list of Table~\ref{tab:playbook-state}: ownership regime, scale and volume, data readiness, regulated share, labor jurisdiction, change capacity and leadership stability, product model, market structure and owner capability.
\end{definition}

The data identify $\lambda$ and the headroom $h(s)$ only through their product $\ell=\lambda h(s)$ (Proposition~\ref{prop:industry}(iv)), and $\alpha$ only within Part~I's calibration range. The position is therefore calibrated as Part~I calibrates its anchor, with one added restriction from Part~I's evidence on spending.

\begin{assumption}[Calibration of a position]\label{ass:playbook-calib}
Growth spend is selling and marketing expense above a maintenance level, $\sigma=\max(s_{\rm SM}-b_M,0)$, where $s_{\rm SM}$ is the selling-and-marketing share of revenue (Section~\ref{sec:data}) and $b_M=\pbMaintBase\%$ is the median share, in the previous year, of firm-years growing at or below the year's median, Part~I's central basis; the rest, $a_{\rm SM}=\min(s_{\rm SM},b_M)$, is a cost line in $a$, and all of $s_{\rm SM}$ and none of it bound $\sigma$. Capital intensity is $k=\kappa_g/g$ for $g>0$, where $\kappa_g$ is capital spending and capitalized software over revenue, and $k=0$ otherwise (the convention of Part~I's anchor); and every firm acts on a net weight with $\alpha(\hat X-k)=\mathsf r_\sigma$, where $\mathsf r_\sigma=\SpendRatioPB$ is the median ratio of growth spend on that basis to excess growth that Section~\ref{sec:empirics} reports (\SpendRatioPBAll{} on all selling and marketing), as in Assumption~\ref{ass:hetero}, under which firms share one weight. Then $\Delta:=\sigma/\mathsf r_\sigma$, the firm's zero-spend growth is $G_i:=g-\Delta$ (Remark~\ref{rem:assumptions}(b)) and $\ell:=\Delta/\sigma^{\alpha}$, with $\ell:=0$ and $G_i=g$ when $\sigma=0$.
\end{assumption}

\begin{lemma}[A calibrated firm is at its own tangency]\label{lem:playbook-calib}
Under Assumption~\ref{ass:playbook-calib}: (a) $\sigma^*(\ell,\hat X)=\sigma$ and $\Delta^*(\ell,\hat X)=\Delta$. (b) For any owner weight $X_{\mathsf o}>k$, $\sigma^*(\ell,X_{\mathsf o})=\sigma[(X_{\mathsf o}-k)/(\hat X-k)]^{1/(1-\alpha)}$ and $\Delta^*(\ell,X_{\mathsf o})=\Delta[(X_{\mathsf o}-k)/(\hat X-k)]^{\alpha/(1-\alpha)}$. (c) For $\alpha\in[0.25,0.45]$, $\hat X-k\in[\mathsf r_\sigma/0.45,\mathsf r_\sigma/0.25]$.
\end{lemma}

The restriction is Part~I's own (one weight for all firms), and its value is class c; for below-median growers it extrapolates a statistic measured on above-median growers, and every lever value that depends on it is reported with that condition. The spend ratio $\mathsf r_\sigma$ depends on the basis (\SpendRatioPB{} netted, \SpendRatioPBAll{} on all selling and marketing) and on Part~I's method of measuring excess growth, and it identifies only the product $\alpha(\hat X-k)$: part (c) is an implication of the calibration, not an estimate of the owner's weight, which Part~I withdraws (the estimate of 7 to 13) and which other baselines put much lower. The playbook therefore uses $\hat X$ in two ways. Rankings that compare an owner's weight with the firm's through the ratio $(X_{\mathsf o}-k)/(\hat X-k)$, as in Proposition~\ref{prop:playbook-rank}(b), and owners whose weight is a multiple of $\hat X-k$ (public, dual class, family) do not depend on the level of $\hat X$. Comparisons between $\hat X$ and a weight computed on its own from rates, horizons and exit prices (a sponsor, a venture investor, a permanent acquirer) do, and they are marked as conditional on that level. % XBOUNDS: level-dependent comparisons are flagged below
The setup enters through the results of Sections~\ref{sec:tasks} to \ref{sec:acquirer}, as Table~\ref{tab:playbook-state} records.

\begin{table}[t]
\centering
\caption{The setup of a firm: what the engine reads, the result each dimension enters through and the grounding of its measure}\label{tab:playbook-state}
\scriptsize
\setlength{\tabcolsep}{3pt}
\begin{tabular}{@{}p{2.6cm}p{4.3cm}p{6.2cm}p{1.3cm}@{}}
\toprule
Dimension & Measure (engine input) & Enters through & Class\\
\midrule
Ownership regime & regime, hold, exit multiple, covenant & weight $X_{\mathsf o}$ (Proposition~\ref{prop:reg-weight}, \ref{prop:acq-pc}(d)); dip budget (\ref{prop:reg-dip}); hold deadline (\ref{prop:reg-speed}(b)); capital hurdle $\zeta^{\rm pe}$ (\ref{prop:reg-value}(d)) & b, d\\
Scale and volume & revenue; process volumes (Tier 3) & $v^*_j$ (\ref{prop:tasks-modes}(e)), $v^{MB}_j$ (\ref{prop:tasks-window}(c)), $v^{BB}$ (\ref{prop:tasks-buildbuy}(a)) & a, d\\
Data readiness & systems inventory (Tier 3) & floor $\underline\iota_j$ (\ref{prop:tasks-review}(d)); $\iota^*_p$ and the big push (\ref{prop:tasks-readiness}(b), (d)); residency cap (\ref{prop:labor-residency}) & d\\
Regulated share & ICFR channels; EU nexus & ITGC lines (\ref{prop:acct-itgc}); cutover (\ref{prop:acct-cutover}); $\mathrm{NPV}^{\rm ICFR}$ (Corollary~\ref{cor:acct-npv}); oversight floor (\ref{cor:ext-floor}) & b, d\\
Labor jurisdiction & sites, headcount (Tier 3) & lead $T^{\rm leg}$ (\ref{prop:labor-lead}); pace (Corollary~\ref{cor:labor-pace}) & a\\
Change capacity, leadership & affected-FTE capacity, decision bandwidth, sponsor tenure & speed limit (\ref{prop:attainable}(c)); pace (\ref{prop:ext-bandwidth}(b)); tenure breakpoint (\ref{prop:ext-shocks}(b)) & b, d\\
Product model & seat or usage tariff (Tier 3) & seat-loss threshold $\nu^*_k$ (\ref{prop:prod-migrate}(b)); pass-through by tariff (\ref{prop:prod-aggregate}(c)) & a, d\\
Market structure & HHI, industry pass-through & durable share $\mathrm{dur}$ in dollars (Corollary~\ref{cor:ext-game-engine}(a)); margin share $1-\bar\varphi^{\rm ind}C_0$ in the score (\ref{prop:tasks-aggregate}(a)); AI-native entry (\ref{prop:ext-should}(b)--(d)) & d\\
Owner capability & platform, reuse, core team & $\mathrm{NPV}^{\mathsf o}_p$ (\ref{prop:acq-owner}); replace or transform (\ref{prop:acq-replace}) & d\\
\bottomrule
\end{tabular}

\smallskip
\begin{minipage}{0.98\linewidth}\scriptsize
Notes: numbers without a prefix are propositions. Classes follow \texttt{notes/assumptions\_register.md}: a sourced, b grounded by precedent, c calibrated from our own data, d scenario value. The engine reads every measure from its input tiers; \texttt{notes/playbook\_engine\_spec.md} lists the fields.
\end{minipage}
\end{table}

\subsection{Lever values and where the ranking switches}\label{sec:playbook-rank}

\begin{definition}[Levers]\label{def:playbook-levers}
A lever moves the position. \emph{Growth-engine automation} saves the share $\delta_\sigma$ of the cost of the growth activities, so that demand pull becomes $\lambda(1-\delta_\sigma)^{-\alpha}$ (Proposition~\ref{prop:tasks-aggregate}(a)); \emph{growth-engine productivity} raises demand pull by the proportion $x_\lambda$ (Corollary~\ref{cor:priorities}(ii)); \emph{overhead} and \emph{engineering} automation cut $a_{\rm GA}$ and $a_{\rm RD}$ by the proportion $x_a$; \emph{cost-of-revenue} automation cuts $c_h$ by the proportion $x_h$; \emph{retention} raises $G_i$ by $dG$; an \emph{AI product} lowers $m_0$ by $D_m$ and raises $G_i$ by $\Delta G$ (Corollary~\ref{cor:prod-dilution}); \emph{owned compute}, or capital for labor generally, removes $da$ and adds $dk$; and an \emph{acquisition} buys revenue at $k^{\rm acq}$ per dollar (Definition~\ref{def:acquired-growth}). For a ranking, the proportional levers take a common proportion $x$.
\end{definition}

Write $\Delta R_\bullet$ for the first-order change in the owner's maximal score $R^*_{X_{\mathsf o}}$ from lever $\bullet$, at the owner's optimum, with spend re-optimized.

\begin{proposition}[Lever values and switch thresholds]\label{prop:playbook-rank}
Let $\bar\sigma=\infty$, $X_{\mathsf o}>k$, and write $\sigma^*$, $\Delta^*$, $g^*=G_i+\Delta^*$ and $\Pi$ for the optimum at $(\ell,X_{\mathsf o})$.
\begin{enumerate}
\item[(a)] (Values.) $\Delta R_\sigma=x\sigma^*$ and $\Delta R_\lambda=x\sigma^*/\alpha$; exactly, $\Pi[(1-x)^{-\alpha/(1-\alpha)}-1]\ge x\sigma^*$ and $\Pi[(1+x)^{1/(1-\alpha)}-1]$. $\Delta R_a=xa$ for $a\in\{a_{\rm GA},a_{\rm RD}\}$, exactly. $\Delta R_h=x(1-\varphi)c_h$, where $\varphi$ is the pass-through elasticity of Lemma~\ref{lem:prices}: for a saving every rival replicates, $\varphi=\bar\varphi^{\rm ind}C_0$, the industry pass-through rate times today's unit cost per dollar of price, with the margin at the lower price (Proposition~\ref{prop:tasks-aggregate}(a)); for a saving no rival can make, the firm's own elasticity. The durable share $\mathrm{dur}=1-\bar\varphi^{\rm ind}(1-\mathrm{HHI})$ of Corollary~\ref{cor:ext-game-engine}(a) is a share of the dollar saving, volume gain included; it values a bundle in dollars and does not enter the score. $\Delta R_G=(X_{\mathsf o}-k)\,dG$, exactly. $\Delta R_k=da-(g^*+\delta_K)dk$. The overhead, engineering, retention and given-demand cost-of-revenue gains raise the saturated score $R_\infty$ one for one; the two growth-engine gains do not, and vanish as $s\to\bar s$.
\item[(b)] (Owner's weight.) The growth engine outranks a line of size $a$ at the same proportion, $\Delta R_\sigma>\Delta R_a$, if and only if
\begin{equation}\label{eq:playbook-xdag}
X_{\mathsf o}>X^\ddagger_0:=k+\frac{a^{1-\alpha}}{\alpha\,\ell},
\end{equation}
equivalently, under Assumption~\ref{ass:playbook-calib}, if and only if $(X_{\mathsf o}-k)/(\hat X-k)>(a/\sigma)^{1-\alpha}$. For productivity, $\Delta R_\lambda>\Delta R_a$ if and only if $X_{\mathsf o}>k+(\alpha a)^{1-\alpha}/(\alpha\ell)$. $X^\ddagger_0$ is the limit of $X^\ddagger$ of \eqref{eq:reg-Xdag} as $\delta_\sigma\to0$ with $da=\delta_\sigma a$. Because $\Delta R_\sigma$, $\Delta R_\lambda$ and $\Delta R_G$ increase in $X_{\mathsf o}$ and $\Delta R_a$ and $\Delta R_h$ do not, a ranking made at the screening weight reverses the owner's whenever $X^{\rm scr}<X^\ddagger_0<X_{\mathsf o}$.
\item[(c)] (Headroom.) $\Delta R_\sigma$ and $\Delta R_\lambda$ are proportional to $h(s)^{1/(1-\alpha)}$, so a line of size $a$ overtakes the growth engine once headroom falls below $h^\ddagger:=h(s)\,(a/\sigma^*)^{1-\alpha}$, and at saturation only the gains that raise $R_\infty$ remain.
\item[(d)] (Cost of revenue.) Cost-of-revenue automation outranks overhead at the same proportion if and only if $(1-\varphi)c_h>a_{\rm GA}$, with $1-\varphi=1-\bar\varphi^{\rm ind}C_0$ for a replicable saving and the firm's own $1-\varphi$ for a saving rivals cannot make (Proposition~\ref{prop:levers-position}(d)). The comparison does not depend on concentration; at the central measured rate of 0.5 (Section~\ref{sec:ext-game}) and a unit cost of 60\% of price, about that of Part~I's anchor, a replicated saving keeps 70\% of its margin gain; if the overhead cut is replicated too, both lines keep the same share and the test is $c_h>a_{\rm GA}$; and under free entry without a scarce input the long-run gain is zero.
\item[(e)] (Retention and AI products.) A point of retention is worth $X_{\mathsf o}-k$ points of margin, and a product plan that dilutes the margin to raise retention pays if and only if $D_m/\Delta G<X_{\mathsf o}-k$; the plan of Corollary~\ref{cor:prod-dilution}'s example, with $D_m/\Delta G=3$, pays at owner weights of 5 and 9 above $k$ and not at 2.5 (illustrative weights). % XBOUNDS
\item[(f)] (Capital.) Owned compute raises the score if and only if $da/dk>g^*+\delta_K$, creates value if and only if $da/dk>r_{\mathsf o}+\delta_K$, and raises a sponsor's value if and only if $da/dk>\zeta^{\rm pe}$. The score hurdle rises with growth. The first two hurdles hold after tax as stated when server capital is fully expensed ($\psi=1$), which is current U.S.\ law for domestic servers acquired after 19 January 2025 unless the firm elects out (Section~\ref{sec:acct-standards}); for servers used predominantly abroad, or after an election out, they are multiplied by $\Upsilon=(1-\psi\tau)/(1-\tau)$, and $\zeta^{\rm pe}$ is before tax.
\item[(g)] (Acquisition.) Acquired revenue raises the score charged with its price if and only if $k^{\rm acq}<X_{\mathsf o}$; under the Rule of 40, if and only if $k^{\rm acq}<1$. At illustrative owner weights of 2.5, 5 and 9, the three prices in the compounder's pro forma statements (Section~\ref{sec:acq-case}), 1.7 (Eventbrite), 2.3 (AOL) and 3.2 (Vimeo) times revenue, all pass at 5 and 9; at 2.5 Vimeo fails, and below 2.3 AOL fails too. % XBOUNDS: the sign flips at X_o = k^acq
\end{enumerate}
\end{proposition}

\noindent\emph{Intuition.} Part~I's envelope results value every lever at the same point, so the ranking is a comparison of line sizes weighted by how each line enters the score. A cost saving is worth one point per point of revenue wherever it falls, but the growth engine's line is the owner's optimal spend $\sigma^*$, which is larger the more the owner values growth and the more headroom remains, while overhead is a fixed line. Part (b) says that the ranking is a comparison of weights: the owner's weight against $X^\ddagger_0$, or, in data the panel identifies, the owner's net weight relative to the firm's revealed net weight against $(a/\sigma)^{1-\alpha}$. A sponsor near entry or a permanent acquirer at its hurdle, whose weight is below the revealed weight unless the firm's zero-spend growth approaches the owner's rate (Proposition~\ref{prop:reg-weight}(c), \ref{prop:acq-pc}(d)), cuts growth spend at once and with it the value of making that spend more productive. Part (c) is the dynamic version: the growth engine is a lever for firms with headroom, and overhead the lever that lasts. Part (d) separates automation that a firm keeps from automation its customers keep, in score terms; the should-versus-could distinction of Section~\ref{sec:ext-competition} makes it in dollars, through the durable share.

\begin{proposition}[Setup gates]\label{prop:playbook-gates}
For an activity, process or program, a lever of Proposition~\ref{prop:playbook-rank} is worth nothing where a gate fails, whatever the position. In the order in which the engine tests them:
\begin{enumerate}
\item[(a)] (Permissibility and physical work.) Licensing, exclusions and data classes remove AI modes, by Lemma~\ref{lem:ext-nested}, and physical activities wait for the world-model lag (Proposition~\ref{prop:ext-feasible}(a)).
\item[(b)] (Capability.) Automation is infeasible before $T^F_j$ (Proposition~\ref{prop:ext-feasible}(a)); under a plateau, activities above it never pass (Proposition~\ref{prop:ext-feasible}(f)).
\item[(c)] (Readiness.) Automation is infeasible below the floor $\underline\iota_j$ (Proposition~\ref{prop:tasks-review}(d)); a process is automated in full only from $\iota^*_p$, and readiness investment pays only as a push to at least $\iota^*_p$ (Proposition~\ref{prop:tasks-readiness}(b), (d)); residency rules cap readiness (Proposition~\ref{prop:labor-residency}(a)).
\item[(d)] (Reliability.) Automation is infeasible at any review share if the catch rate is below $1-\bar\varepsilon_j/e_j$ on the catchable rate (Propositions~\ref{prop:tasks-review}(d), \ref{prop:tasks-halluc}(b)), higher still under automation bias (Proposition~\ref{prop:tasks-bias}(d)); augmentation remains (Proposition~\ref{prop:tasks-augment}).
\item[(e)] (Cost and scale.) Automation is valued against the cheapest alternative (Proposition~\ref{prop:tasks-modes}(c)); below $v^*_j$, $v^{MB}_j$ or $v^{BB}$ the firm should keep, outsource or buy (Propositions~\ref{prop:tasks-modes}(e), \ref{prop:tasks-window}(c), \ref{prop:tasks-buildbuy}(a)); run costs set a floor no price removes (Proposition~\ref{prop:tasks-fullcost}(b)); a higher owner rate raises every threshold (Proposition~\ref{prop:reg-speed}(a)).
\item[(f)] (Transition.) A transition pays only if its risk-adjusted NPV, with the ICFR terms for a financial process, is positive (Proposition~\ref{prop:tasks-transition}(b), Corollary~\ref{cor:acct-npv}); a pilot pays only inside its band (Proposition~\ref{prop:tasks-pilot}(c)); a program whose embedding window exceeds $1/r_D$ loses value to leadership and strategy hazards (Proposition~\ref{prop:ext-shocks}(b)).
\item[(g)] (Pace.) A program lasts at least the largest of its change-capacity, bandwidth and talent bounds (Corollary~\ref{cor:labor-pace}); an owner's dip budget and hold deadline cap what can be done in time (Propositions~\ref{prop:reg-dip}(b), \ref{prop:reg-speed}(b)); dismissals wait for the legal lead (Proposition~\ref{prop:labor-lead}(a)).
\item[(h)] (Competition and the owner.) A replicable saving keeps only its durable share, so its automation can be defensive (Corollary~\ref{cor:ext-game-engine}(b)); after AI-native entry binds the incumbent keeps its savings but survives only if its switching premium and scarce inputs cover its readiness gap (Proposition~\ref{prop:ext-should}(b), (c)); and a transition that loses money for a standalone firm can pay for an owner with platform capability (Proposition~\ref{prop:acq-owner}(c)).
\end{enumerate}
\end{proposition}

\begin{corollary}[Levers that do not pay]\label{cor:playbook-nopay}
By Propositions~\ref{prop:playbook-rank} and \ref{prop:playbook-gates}, and with the results cited there:
\begin{enumerate}
\item[(i)] growth-engine levers, for an owner whose weight is at most $X^\ddagger_0$ and for a firm whose headroom is below $h^\ddagger$;
\item[(ii)] cost-of-revenue automation as a margin lever where $\mathrm{dur}\,c_h<a_{\rm GA}$: it remains a defensive lever, to be done because rivals will;
\item[(iii)] owned compute wherever the operating cost it saves per dollar of capital is below $g^*+\delta_K$, which excludes it first for the fastest growers;
\item[(iv)] automation of steps that fail the reliability gate at measured catch rates, which is where regulated finance steps sit at the engine's defaults, and readiness investment that stops short of $\iota^*_p$;
\item[(v)] in-house builds below the build-or-buy and make-or-buy volumes, and in-house automation after an AI-enabled vendor has closed the window;
\item[(vi)] AI products whose dilution ratio exceeds $X_{\mathsf o}-k$, a test that rejects at a screening weight plans the firm's own weight accepts, and whose outcome depends on the level of that weight; % XBOUNDS
\item[(vii)] replacing an organization instead of transforming it when the owner's weight is high or customers are sensitive, since the customer cost of replacement is proportional to $X_{\mathsf o}$ (Proposition~\ref{prop:acq-replace}(b)), a level-dependent test; % XBOUNDS
\item[(viii)] acquisitions priced above $X_{\mathsf o}$ per dollar of revenue, a test that depends on the level of the owner's weight (Proposition~\ref{prop:playbook-rank}(g)), and, for a base that shrinks, growth bought only to stand still (Proposition~\ref{prop:acq-limits}(a)); % XBOUNDS
\item[(ix)] programs that do not embed within $1/r_D$, unless split into wins that do.
\end{enumerate}
\end{corollary}

\subsection{Archetypes}\label{sec:playbook-archetypes}

\begin{definition}[Archetypes]\label{def:playbook-archetypes}
Let the population be the \pbCutNpop{} firm-years of Part~I's core sample in fiscal 2023--25 with a free-cash-flow margin after stock-based compensation (estimand E1), and the members of each cell its \pbCutNEtwo{} non-acquisition firm-years (estimand E2). With population quantiles of growth $g$ ($Q_{25}=\pbCutGtwentyfive$, median $\pbCutGmed$, $Q_{75}=\pbCutGseventyfive$), margin $m$ ($\pbCutMtwentyfive$, $\pbCutMmed$, $\pbCutMseventyfive$) and score $R=g+m$ ($Q_{75}=\pbCutRseventyfive$, $Q_{95}=\pbCutRninetyfive$), all in percent: A1, the \emph{growth corner}, has $g\ge Q_{75}$ and $m<Q_{25}$; A2, the \emph{balanced upper quartile}, has $Q_{75}\le R<Q_{95}$ with $g$ and $m$ at or above their medians; A3, the \emph{mature cash generator}, has $g$ below the median and $m\ge Q_{75}$; A4, the \emph{sub-40 grower}, has $g$ between the median and $Q_{75}$ and $m$ below the median; A5, the \emph{stalled} firm, has $g$ and $m$ below their medians; A6, the \emph{near-frontier elite}, has $R\ge Q_{95}$, Part~I's calibration anchor, and takes precedence over the other cells. A7, the \emph{acquisition-driven compounder}, is the serial acquirer of Section~\ref{sec:acq-case}, with reported and organic growth from its filings. Each archetype's position and cost lines are its cell's medians.
\end{definition}

Only \pbCutNpassforty{} non-elite, non-acquisition firm-years of 2023--25 pass the Rule of 40 with both components above their medians, which is why A2 is defined by the score's upper quartile rather than by the pass mark. Table~\ref{tab:playbook-archetypes} reports the positions and their calibrated states. The near-frontier elite has the highest zero-spend growth ($G_i=\pbASixGi\%$, all of its growth, since its S\&M is below the maintenance base), the growth corner the largest growth spend (S\&M of $\pbAOneSM\%$ of revenue, \pbAOneSigma{} points above the base) and the stalled firm a base that shrinks without spend ($G_i=\pbAFiveGi\%$). For the compounder, reported growth was \pbASevenGrep\% in the second quarter of 2026 and organic growth \pbASevenGorg\% on the company's definition, for scores of \pbASevenRrep{} and \pbASevenRorg{} points on its operating margin; its net revenue retention implies $G_i=\pbASevenGi\%$.

\IfFileExists{part2/tables/playbook_archetypes.tex}{\begin{table}[t]
\centering
\caption{Archetypes: cell medians of 2023--25 non-acquisition core firm-years and their calibrated Part~I states (Assumption~\ref{ass:playbook-calib}, $\alpha$ the engine's default)}\label{tab:playbook-archetypes}
\scriptsize
\setlength{\tabcolsep}{3pt}
\resizebox{\textwidth}{!}{\begin{tabular}{@{}lrrrrrrrrrrrl@{}}
\toprule
 & $n$ (firms) & $g$ & $m$ & S\&M & $\sigma$ & G\&A & R\&D & GM & $k$ & $G_i$ & $\hat X$ & Ranks $g/m/R$\\
\midrule
A1 Growth corner & 28 (23) & 29.1 & \ensuremath{-}23.0 & 50 & 25 & 24 & 32 & 70 & 0.09 & 4.1 & 3.4 & p88/p12/p32 \\
A2 Balanced upper quartile & 43 (30) & 20.9 & 13.8 & 26 & 1 & 10 & 20 & 79 & 0.10 & 19.4 & 3.4 & p73/p79/p86 \\
A3 Mature cash generator & 52 (29) & 7.7 & 19.7 & 19 & 0 & 10 & 15 & 80 & 0.19 & 7.7 & 3.5 & p31/p88/p73 \\
A4 Sub-40 grower & 48 (35) & 17.1 & \ensuremath{-}5.4 & 43 & 18 & 19 & 26 & 74 & 0.12 & \ensuremath{-}1.0 & 3.4 & p64/p32/p38 \\
A5 Stalled & 84 (49) & 2.1 & \ensuremath{-}12.2 & 31 & 6 & 20 & 25 & 70 & 1.21 & \ensuremath{-}4.2 & 4.5 & p18/p22/p14 \\
A6 Near-frontier elite & 13 (7) & 36.3 & 22.4 & 20 & 0 & 15 & 21 & 74 & 0.04 & 36.3 & 3.3 & p95/p92/p98 \\
A7 Acquisition-driven compounder & 1 & 95$^\dagger$ & \ensuremath{-}120$^\dagger$ & 10 & 0 & 25 & 9 & 66 & -- & \ensuremath{-}5.1 & 3.3 & F-1, 6-K \\
\bottomrule
\end{tabular}}

\smallskip
\begin{minipage}{0.98\linewidth}\scriptsize
Notes: percent of revenue except $n$, $k$ and $\hat X$; $m$ is the free-cash-flow margin after stock-based compensation; GM gross margin. $\sigma$ is growth spend on the maintenance-netted basis, S\&M less a maintenance base of 25.1\% of revenue (zero below it; Part~I's central basis), the rest of S\&M being a cost line; $\mathsf r_\sigma=0.99$, the median ratio of S\&M above the base to excess growth. $k$ is capital spending and capitalized software over revenue divided by $g$; $G_i=g-\sigma/\mathsf r_\sigma$; $\hat X=k+\mathsf r_\sigma/\alpha$. $^\dagger$A7: fiscal 2025 reported growth and the fully charged cash margin (operating cash flow plus investing cash flow, which includes acquisitions) from the F-1; its $G_i$ is $\ln$ of its 2025 net revenue retention; its S\&M is below the base, so it has no growth spend on this basis and $\hat X$ is the common weight. Generated by \texttt{code/playbook/run\_archetypes.py}.
\end{minipage}
\end{table}
}{}

\subsection{The playbook table}\label{sec:playbook-table}

Table~\ref{tab:playbook} is the playbook at a reference setup: a public owner, today's readiness 0.4 in every function (the low-readiness case of Section~\ref{sec:tasks}), a standalone firm and the engine's default market concentration. Each cell traces to a result. The binding constraint has two parts: the cost line in which the archetype's median falls furthest short of the top decile of operating margin in its growth band, on the bands and sample of Part~I's descriptive decomposition of the margin gap (growth bands with edges at 0, 10, 20, 30, 40 and 60\%; \texttt{code/decompose\_gap.py}), and the capacity the engine's program uses most in its first year. The levers are ranked by Proposition~\ref{prop:playbook-rank}(a) at the public owner's weight, and the levers to avoid follow Corollary~\ref{cor:playbook-nopay}. Because the panel's firms are public, the weight that rationalizes an archetype's spending is a public owner's net weight under its earnings pressure, $X^{\rm my}$ of \eqref{eq:reg-Xmy}; we therefore take the public owner's weight to be $\hat X$ and recover the same firm's weight without the pressure by inverting \eqref{eq:reg-Xmy} at the public preset of Table~\ref{tab:reg-presets}, the starting point for the dual-class and family owners of Table~\ref{tab:playbook-grid}. The first wave is the engine's: \texttt{code/playbook/run\_archetypes.py} writes each archetype's cell medians as Tier~0 financials, sets readiness through a Tier~3 systems inventory and runs the engine, which builds the activity map, prices every mode, optimizes the program under the public regime and orders it into waves. Wave~0, measurement, always comes first.

\begin{table}[p]
\centering
\caption{The playbook by archetype at the reference setup (public owner, today's readiness 0.4, standalone firm). Lever values are points of the maximal score per 10\% of a line automated, at the owner's weight (Proposition~\ref{prop:playbook-rank}(a)); levers to avoid cite the items of Corollary~\ref{cor:playbook-nopay}; the first wave is the engine's (Section~\ref{sec:playbook-table})}\label{tab:playbook}
\scriptsize
\setlength{\tabcolsep}{3pt}
\renewcommand{\arraystretch}{1.1}
\newcolumntype{Q}[1]{>{\raggedright\arraybackslash}p{#1}}
\begin{tabular}{@{}Q{1.8cm}Q{2.1cm}Q{1.9cm}Q{2.4cm}Q{4.0cm}Q{2.9cm}@{}}
\toprule
Archetype & Position (panel) & Binding: largest gap; scarcest capacity & First two levers (points) & Levers to avoid & First wave (engine)\\
\midrule
A1 Growth corner & $g=30\%$ (p90), $m=-23\%$ (p13); $n=28$ & S\&M 27 pts; bandwidth & 1. R\&D 3.2; 2. growth engine 2.4 & (iii) owned compute below 0.55 of opex per \$ of capital; (ii) cost of revenue as a margin lever (keeps 0.53 of the cut); (i) growth engine under a sponsor ($X=0.3<3.9$); (iv) automating 27 ICFR steps that fail the catch-rate gate at $\iota=0.9$ & Lead to close (2026; NPV \$361M) \\
\addlinespace
A2 Balanced upper quartile & $g=21\%$ (p72), $m=14\%$ (p79); $n=42$ & R\&D 4 pts; bandwidth & 1. R\&D 2.0; 2. G\&A 1.0 & (iii) owned compute below 0.46 of opex per \$ of capital; (ii) cost of revenue as a margin lever (keeps 0.59 of the cut); (i) growth engine: no S\&M above the maintenance base; (iv) automating 29 ICFR steps that fail the catch-rate gate at $\iota=0.9$ & Lead to close (2026; NPV \$364M) \\
\addlinespace
A3 Mature cash generator & $g=8\%$ (p31), $m=19\%$ (p87); $n=50$ & R\&D 4 pts; bandwidth & 1. R\&D 1.5; 2. G\&A 1.1 & (iii) owned compute below 0.33 of opex per \$ of capital; (ii) cost of revenue as a margin lever (keeps 0.62 of the cut); (i) growth engine: no S\&M above the maintenance base; (iv) automating 29 ICFR steps that fail the catch-rate gate at $\iota=0.9$ & Ticket to resolution (2026; NPV \$53M) \\
\addlinespace
A4 Sub-40 grower & $g=17\%$ (p64), $m=-5\%$ (p34); $n=48$ & S\&M 27 pts; bandwidth & 1. R\&D 2.6; 2. G\&A 1.9 & (iii) owned compute below 0.42 of opex per \$ of capital; (ii) cost of revenue as a margin lever (keeps 0.58 of the cut); (i) growth engine under a sponsor ($X=1.0<4.5$); (iv) automating 28 ICFR steps that fail the catch-rate gate at $\iota=0.9$ & Product and engineering (other work), Finance and accounting (other work) (2026; NPV \$89M) \\
\addlinespace
A5 Stalled & $g=2\%$ (p17), $m=-13\%$ (p22); $n=80$ & R\&D 15 pts; bandwidth & 1. R\&D 2.5; 2. G\&A 2.0 & (iii) owned compute below 0.27 of opex per \$ of capital; (ii) cost of revenue as a margin lever (keeps 0.49 of the cut); (i) growth engine under a sponsor ($X=-0.1<9.3$); (iv) automating 28 ICFR steps that fail the catch-rate gate at $\iota=0.9$ & Professional services (other work) (2026; NPV \$7M) \\
\addlinespace
A6 Near-frontier elite & $g=35\%$ (p95), $m=22\%$ (p92); $n=14$ & R\&D 6 pts; bandwidth & 1. R\&D 2.0; 2. G\&A 1.6 & (iii) owned compute below 0.60 of opex per \$ of capital; (ii) cost of revenue as a margin lever (keeps 0.63 of the cut); (i) growth engine: no S\&M above the maintenance base; (iv) automating 28 ICFR steps that fail the catch-rate gate at $\iota=0.9$ & Lead to close, Product and engineering (other work) (2026; NPV \$422M) \\
\addlinespace
A7 Acquisition-driven compounder & reported $g=126\%$, organic 3\%; $m_{op}=34\%$ (Q2 2026) & base shrinks ($G_i=-5.1\%$): the treadmill & 1. owner capability (platform, rebuild); 2. G\&A and R\&D cost-out & (viii) acquisitions above $X^G=0.77$ times revenue; (i) growth spend & \emph{not generated: the engine models no acquisitions or platform owner (W12)} \\
\bottomrule
\end{tabular}

\smallskip
\begin{minipage}{0.98\linewidth}\scriptsize
Notes: positions are cell medians of 2023--25 non-acquisition core firm-years of the audited v3 sample (Definition~\ref{def:playbook-archetypes}); p$q$ is the population quantile rank. Binding: the cost line with the largest gap of the archetype's median to the top decile of operating margin in its growth band (bands with edges at 0, 10, 20, 30, 40 and 60\% and the sample of \texttt{code/decompose\_gap.py}), in points of revenue, and the engine's most used capacity in the first plan year. Lever values at the public owner's weight, the revealed weight $\hat X$ (the panel's firms are public, so their spending reveals a public owner's net weight $X^{\rm my}$ of eq.~\eqref{eq:reg-Xmy}); the growth engine is automation of the selling and marketing that buys growth, S\&M above a maintenance base of 25.7\% of revenue (first order $x\sigma^*$; Part~I's central, maintenance-netted basis, with S\&M up to the base a cost line); R\&D treats all R\&D as maintenance, an upper bound; cost of revenue at the margin share a replicated cut keeps, $1-\bar\varphi^{\rm ind}C_0$, at the engine's central pass-through rate of 0.50 and the archetype's unit cost per dollar of price $C_0$ (Part~I's round-2 mapping: the rate applies to the unit-cost saving and the margin is computed at the lower price; the durable share 0.625, a dollar share, does not enter the score). With all S\&M as growth spend, Part~I's upper bound, the growth engine would be the first lever at six archetypes, and with none of it at none (Section~\ref{sec:playbook-limits}). Levers to avoid: owned compute below its score hurdle $g^*+\delta_K$ (after tax with domestic servers fully expensed, $\psi=1$; times $(1-\psi\tau)/(1-\tau)$ for servers abroad); growth engine under a sponsor when the sponsor's weight is below $X^\ddagger_0$ for the line that ranks first; ICFR steps counted at readiness 0.9. Waves from the engine's roadmap; residual work streams are named by function. Generated by \texttt{code/playbook/run\_archetypes.py}.
\end{minipage}
\end{table}


Five readings follow, on Part~I's central, maintenance-netted basis for growth spend (Section~\ref{sec:playbook-limits} gives the all-S\&M and no-S\&M bounds). First, whether the growth engine or overhead comes first depends on how much selling and marketing buys growth. At the public owner's weight, cutting general and administrative cost comes before automating the growth engine at \pbNGAAheadGENet{} of the six panel archetypes. At that weight, which is the revealed weight, the growth engine outranks overhead exactly when growth spend exceeds the overhead line (Proposition~\ref{prop:playbook-rank}(b)), so it leads only where S\&M is well above the maintenance level: at the growth corner (\pbGEAheadNetList), whose growth spend of \pbAOneSigma\% of revenue exceeds its overhead of \pbAOneGA\%. At the sub-40 grower, also well above the base, growth spend of \pbAFourSigma\% falls short of overhead of \pbAFourGA\%, and at \pbSMNetZeroList{} S\&M is below the base, so there is no growth spend to automate. Engineering is the first lever at \pbNEngFirstRevNet{} of the six, on a value that treats all R\&D as maintenance and is therefore an upper bound. With all S\&M as growth spend the growth engine would be first at \pbNGEFirstRevAll{} of the six, and with none of it at \pbNGEFirstRevNone.
Second, the owner moves the ranking only where there is growth spend to reweight. Every switch in Table~\ref{tab:playbook-grid} is a crossing of $X^\ddagger_0$ in \eqref{eq:playbook-xdag}, against engineering or overhead. A family owner, whose weight is higher, puts the growth engine first at \pbNGEFirstFam{} archetype (\pbGEFirstFamList), and a dual-class owner at \pbNGEFirstHyb. Under a sponsor at the entry of a six-year hold the growth engine is first at \pbNGEFirstPE{} archetype, and under a permanent acquirer at its unlevered hurdle at \pbNGEFirstPC{} of the \pbNDefinedPC{} archetypes at which its going-concern weight is defined; both cut growth spend where there is any (both comparisons are conditional on the level of $\hat X$, which the calibration sets at $\hat X-k=\mathsf r_\sigma/\alpha$; a lower level turns more of them toward the growth engine); % XBOUNDS
the acquirer's weight is not defined for \pbPCUndefined, whose zero-spend growth exceeds the hurdle of \pbPCRate\% (Table~\ref{tab:playbook-grid}). At the elite no owner changes the ranking: its calibrated zero-spend growth of \pbASixGi\% exceeds the sponsor's rate of \pbPERate\%, so \eqref{eq:reg-Xexit} gives the sponsor a weight of \pbASixXPE, above the revealed \pbASixXhat, but on this basis it has no growth spend for that weight to act on; this rests on the class-c calibration of $G_i$ and on the level of $\hat X$. % XBOUNDS
Third, for slow growers the gap evidence and the lever values point to the same line: their margins fall short of their band's top decile mostly in \pbAThreeGap{} (A3) and \pbAFiveGap{} (A5), and engineering is their first lever. The agreement is weaker than it looks. The gap measures where better firms differ, which can reflect product mix and is not a lever (Section~\ref{sec:empirics}), and the engineering value is an upper bound. Fourth, owned compute is to be avoided at every panel position unless the operating cost it saves per dollar of capital exceeds a hurdle that ranges from \pbAFivePubCaph{} for the stalled firm to \pbASixPubCaph{} for the elite (Proposition~\ref{prop:playbook-rank}(f), after tax with domestic servers fully expensed), and cost-of-revenue automation is among the first two levers for \pbNCoRTopTwo{} of the archetypes: at the engine's central pass-through rate of \pbPtRate, a replicated cut keeps \pbKeepMin{} to \pbKeepMax{} of its margin gain at the archetypes' unit costs (part (d)), whatever the concentration. Fifth, what limits the engine's program in its first year is the same for every archetype: decision bandwidth, the most used capacity for \pbNBindBandwidth{} of them, with load at \pbBandwidthUseMin{} to \pbBandwidthUseMax{} times its breakpoint (Proposition~\ref{prop:ext-bandwidth}(b)), so the first wave is small whatever the position.

The compounder's row cannot yet be generated by the engine, which models neither acquisitions nor a platform owner's capability (\texttt{notes/playbook\_engine\_spec.md}, item W12). Its position-level entries come from Proposition~\ref{prop:playbook-rank}(g) and Section~\ref{sec:acquirer}: at its unlevered hurdle and the margin a permanent dollar earns, its growth weight is $X^G=\pbASevenXpc$, so every reported acquisition price lowers the charged score at its own weight (a weight computed from its hurdle and margin, not from $\hat X$); AOL, at \pbASevenAOLprice{} times revenue, changes it by \pbASevenAOL{} points per point of revenue acquired; acquisitions pay it only through what its capability adds to the business (Proposition~\ref{prop:acq-pc}(b)), which is the replace-and-rebuild path of Proposition~\ref{prop:acq-replace}. Reorganization expense, the cost of that path, was \pbASevenReorg\% of revenue in 2025.

\subsection{From levers to steps: the executable playbook}\label{sec:playbook-steps}

The archetype table is a strategy; the engine also produces the work plan. Decisions are made at the level of the engine's activity unit, an O*NET detailed work activity (DWA) performed in a process step or in a function's residual work, and executed at the level of the process step, where workflows, owners, systems, controls and waves attach (Section~\ref{sec:tasks-setup}). O*NET task statements linked to each DWA are attached as descriptors and evidence keys, not as decision units. For every archetype, readiness case and step the runner writes one record with eleven groups of fields (\texttt{notes/playbook\_engine\_spec.md}, Section~4): the mode now and at target readiness, for the activity and for the step within its process's best plan; the binding constraint on automation, in the order of Proposition~\ref{prop:playbook-gates}; preconditions; go and no-go gates (the minimum catch rate, the tolerated residual error, the minimum volume and the pilot band); the wave and the steps that move with it; roles and FTE effects; KPIs; controls (ICFR status, channel and condition, ITGC and interface lines); evidence keys with their grounding class, including the Anthropic Economic Index usage and automation share and the capability date from the horizon fit; a confidence tier; and switch conditions with the firm's distance to each threshold.

\begin{definition}[Confidence tier]\label{def:playbook-tier}
A recommendation is in Tier~C if the decision to automate or not, or to use an AI mode or not, changes at an end of the stated range of some class-d input with the others at their defaults, or if the share of the engine's Monte Carlo draws in which its process's value is positive is below 0.8, a setting; the inputs are named. Otherwise it is in Tier~A if it is bound by a feasibility filter that the backtest of Section~\ref{sec:backtest} supports (capability horizon, physical work, permissibility), with sourced costs, and in Tier~B if it rests on the cost crossover or a reliability or readiness gate with sourced inputs, which the backtest does not validate.
\end{definition}

Every step of the process map has a record, and the activity records cover the work a software firm's functions perform. Of the \pbCovSteps{} steps, \pbCovStepsPop{} are in every archetype's activity map; the others belong to activities the engine's priors do not staff in a software firm and are kept with a flag. The activity records cover \pbCovUnits{} activity units per configuration and \pbCovDWA{} distinct DWAs, the ones a software firm's functions perform, out of \pbCovDWAAll. Across archetypes and readiness cases, \pbTierStepA\% of step decisions are in Tier~A, \pbTierStepB\% in Tier~B and \pbTierStepC\% in Tier~C. Table~\ref{tab:playbook-functions} summarizes the activity decisions by function for archetype A2; Table~\ref{tab:playbook-processes} in Appendix~\ref{app:playbook-tables} gives the step decisions by process, and Table~\ref{tab:playbook-steps} lists every step.

\IfFileExists{part2/tables/playbook_functions.tex}{\begin{table}[t]
\centering
\caption{Activity decisions by function for archetype A2 (shares of in-house FTE): automated and augmented now, at readiness 0.4 and 0.9, the most common binding constraint at 0.4, and the confidence tiers}\label{tab:playbook-functions}
\scriptsize
\setlength{\tabcolsep}{3pt}
\begin{tabular}{@{}lrrrrrp{4.2cm}r@{}}
\toprule
 & & \multicolumn{2}{c}{$\iota=0.4$ (\%)} & \multicolumn{2}{c}{$\iota=0.9$ (\%)} & & Tier A/B/C\\
\cmidrule(lr){3-4}\cmidrule(lr){5-6}
Function & FTE & A & U & A & U & Binding constraint (most FTE) & (\% FTE)\\
\midrule
Executive and strategy & 149 & 0 & 86 & 0 & 86 & feasibility: capability horizon & 98/2/1 \\
Product and engineering & 2543 & 0 & 83 & 0 & 83 & feasibility: capability horizon & 98/0/2 \\
Marketing & 727 & 0 & 84 & 0 & 79 & feasibility: capability horizon & 77/5/18 \\
Sales & 2505 & 0 & 80 & 0 & 83 & feasibility: capability horizon & 70/30/0 \\
Customer support & 1571 & 0 & 37 & 4 & 29 & feasibility: capability horizon & 46/5/49 \\
Customer success and account management & 739 & 0 & 63 & 0 & 68 & feasibility: capability horizon & 59/12/29 \\
Professional services & 916 & 0 & 77 & 0 & 78 & feasibility: capability horizon & 90/8/2 \\
Service delivery and operations & 36 & 0 & 62 & 0 & 65 & feasibility: capability horizon & 71/5/24 \\
Supply chain, manufacturing and logistics & 42 & 0 & 59 & 0 & 47 & feasibility: capability horizon & 64/22/14 \\
Procurement & 39 & 0 & 62 & 0 & 48 & feasibility: capability horizon & 38/19/43 \\
Finance and accounting & 467 & 0 & 50 & 0 & 42 & feasibility: capability horizon & 46/6/48 \\
Human resources & 176 & 0 & 41 & 0 & 41 & feasibility: capability horizon & 45/10/45 \\
Legal & 83 & 0 & 38 & 0 & 38 & permissibility & 30/0/69 \\
Risk, compliance, audit and governance & 79 & 0 & 51 & 0 & 64 & feasibility: capability horizon & 61/2/38 \\
Information technology and security & 111 & 0 & 52 & 0 & 59 & feasibility: capability horizon & 67/0/33 \\
Technical operations and hosting & 555 & 0 & 72 & 0 & 78 & feasibility: capability horizon & 98/0/1 \\
Facilities and real estate & 81 & 0 & 27 & 0 & 28 & feasibility: physical & 86/6/8 \\
Business operations and transformation & 65 & 0 & 73 & 0 & 70 & feasibility: capability horizon & 59/7/34 \\
Corporate affairs and investor relations & 44 & 0 & 69 & 0 & 62 & feasibility: capability horizon & 78/3/18 \\
\bottomrule
\end{tabular}

\smallskip
\begin{minipage}{0.98\linewidth}\scriptsize
Notes: activity decisions of the engine (cheapest feasible mode by its choice cost, 2026); tiers of Definition~\ref{def:playbook-tier}. Generated by \texttt{code/playbook/run\_archetypes.py}; rows in \texttt{data/playbook/playbook\_activities.csv}.
\end{minipage}
\end{table}
}{}

Two findings concern the evidence rather than the firms. First, attributes measured at the DWA level travel badly to steps. The physical-task share of a DWA averages over every occupation that performs it, so ``Reconcile records of sales or other financial transactions'', which cashiers also perform, carries a physical share above one half, and the engine treats a software firm's reconciliations as physical work that waits for the world-model lag: in the reference runs \pbStepsPhysicalFinance{} procurement and finance steps are bound by the physical flag in every archetype (\pbStepsPhysicalFinanceList), and the invoice chain's matching step (\pbStepsPhysicalFinanceSome) in one. The playbook needs step-level overrides of these attributes, which the engine specification requests (item W5); until they exist these rows carry Tier~A for a constraint that does not apply to them, and the reader should treat them as unresolved. Second, where the physical flag does not bind, readiness binds first and reliability next. At readiness 0.4, \pbICFRReadShare\% of in-scope ICFR step decisions are bound by the readiness floor and \pbICFRRelShareNow\% by reliability; at 0.9, \pbICFRRelShare\% are bound by the catch-rate gate of Proposition~\ref{prop:playbook-gates}(d). The matching step shows the sequence. For A2 it is outsourced today and bound by readiness (effective readiness \pbInvTwoLowRead{} against a floor of \pbInvTwoLowFloor); at readiness 0.9 it is bound by the catch rate, \pbInvTwoHighDnow{} against the \pbInvTwoHighDstar{} that the tolerated residual error of \pbInvGate\% requires, as in Section~\ref{sec:tasks}; for the stalled firm the gate is \pbInvFiveHighDstar{} against \pbInvFiveHighDnow. At the engine's defaults the finance playbook is integration first, and even at readiness 0.9 the activity decision for in-scope ICFR steps is mostly to outsource (\pbICFROutHigh\%) or keep the work in house (\pbICFRHouseHigh\%); augmentation is the decision for \pbICFRAugHigh\% and automation for \pbICFRAutoHigh\%.

\subsection{Limits}\label{sec:playbook-limits}

The ranking rests on first-order envelope values at the owner's optimum; exact values are reported where they differ, and the growth-engine values are lower bounds. The position calibration assumes one net weight for all firms, a class-c value measured on above-median growers whose level the data do not identify, and reads it as a public owner's net weight. It counts as growth spend only S\&M above a maintenance base of \pbMaintBase\% of revenue, Part~I's central basis, with all S\&M and none of it as bounds, and the lead finding depends on that choice. With all S\&M as growth spend, Part~I's upper bound, and a spend ratio of \SpendRatioPBAll, the growth engine outranks general and administrative cost at the revealed weight at \pbNGEAheadGAAll{} of the six positions, against \pbNGEAheadGANet{} on the netted basis, and is the first of the four levers at \pbNGEFirstRevAll; with none of it, the growth-engine levers vanish and engineering is first at \pbNEngFirstRevNone. The basis decides how a saving in S\&M is valued, not whether it is large: on the netted basis S\&M up to the base is a cost line, larger than overhead at \pbNSMMaintAboveGA{} of the six positions, so automating the same share of all S\&M outranks overhead at \pbNSMAutoAheadGA{} on every basis, as a durable cost cut for the maintenance part and as a growth-engine gain, which depends on the owner's weight and the headroom, for the rest. It also treats all R\&D as maintenance, an upper bound on the engineering lever. The engine's own regime runs use the same public, dual-class, sponsor and acquirer weights (specification items W1, W14, W15 and W17) but value whole programs with their own presets, so Table~\ref{tab:playbook-grid} reports their values beside the lever map rather than inside it. Archetype positions are cell medians, not firms. The cost-of-revenue values rest on the scenario pass-through rate, the owner weights on regime presets that are mostly surveys and scenario values (Table~\ref{tab:reg-presets}), and every step-level number on the engine's defaults, with the tier saying which recommendations would change within their ranges. The compounder's row is position-only until the engine models owner capability.
